\documentclass{article}

\usepackage{amsmath,amssymb}
\usepackage{xurl}
\usepackage[hidelinks]{hyperref}
\setlength{\emergencystretch}{2em}

% Shared commands for notes in docs/paper-gaps/.

\DeclareUnicodeCharacter{2080}{\ifmmode{}_{0}\else\textsubscript{0}\fi}
\DeclareUnicodeCharacter{2081}{\ifmmode{}_{1}\else\textsubscript{1}\fi}
\DeclareUnicodeCharacter{2082}{\ifmmode{}_{2}\else\textsubscript{2}\fi}
\DeclareUnicodeCharacter{2099}{\ifmmode{}_{n}\else\textsubscript{n}\fi}

\newcommand{\Lean}{\textsc{Lean}}
\ExplSyntaxOn
\NewDocumentCommand{\leanid}{m}{
  \ifmmode
    \textsf{\detokenize{#1}}
  \else
    \group_begin:
    \urlstyle{sf}
    \tl_set:Nn \l_tmpa_tl {#1}
    \tl_replace_all:Nnn \l_tmpa_tl {\_} {_}
    \exp_args:NV \path \l_tmpa_tl
    \group_end:
  \fi
}
\ExplSyntaxOff
\providecommand{\tr}{\operatorname{tr}}
\newcommand{\tnleanIssueBase}{https://github.com/LionSR/TNLean/issues/}
\newcommand{\tnleanPrBase}{https://github.com/LionSR/TNLean/pull/}
\newcommand{\tnleanDocBase}{https://sirui-lu.com/TNLean/docs/}
\newcommand{\ghissue}[1]{\href{\tnleanIssueBase#1}{GitHub issue \##1}}
\newcommand{\ghpr}[1]{\href{\tnleanPrBase#1}{PR \##1}}
\newcommand{\doclink}[1]{\href{\tnleanDocBase#1.html}{\path{#1}}}

% Dirac notation and the identity operator, as in blueprint/src/macros/common.tex.
% \providecommand keeps note-local definitions authoritative.
\usepackage{dsfont}
\providecommand{\ket}[1]{|#1\rangle}
\providecommand{\bra}[1]{\langle #1|}
\providecommand{\braket}[2]{\langle #1 | #2 \rangle}
\providecommand{\ketbra}[2]{|#1\rangle\!\langle #2|}
\providecommand{\Id}{\mathds{1}}
\providecommand{\one}{\mathds{1}}
\providecommand{\C}{\mathbb{C}}
\providecommand{\MN}[1]{M_{#1}(\C)}
\providecommand{\MD}{M_D(\C)}

% External referencing for paper sources.  The build helper writes sanitized
% auxiliary files under build/paper-aux/ and excludes duplicated source labels.
\usepackage{xr}

% Import a generated paper auxiliary file with a caller-supplied label prefix.
% The candidate paths support builds from docs/paper-gaps/, the repository root,
% and build/paper-aux/ itself.
\newcommand{\importpaperaux}[2]{%
  \IfFileExists{../../build/paper-aux/#2.aux}{%
    \externaldocument[#1]{../../build/paper-aux/#2}%
  }{%
    \IfFileExists{build/paper-aux/#2.aux}{%
      \externaldocument[#1]{build/paper-aux/#2}%
    }{%
      \IfFileExists{#2.aux}{%
        \externaldocument[#1]{#2}%
      }{}%
    }%
  }%
}

% General external reference macros
\newcommand{\paperref}[2]{\ref{#1-#2}}
\newcommand{\papereqref}[2]{(\ref{#1-#2})}

% CPSV16 source (1606.00608 / MPDO-22-12-17-2.tex) external document and shortcuts
\importpaperaux{cpsv16-}{1606.00608/MPDO-22-12-17-2-tnlean}
\newcommand{\cpsvref}[1]{\paperref{cpsv16}{#1}}
\newcommand{\cpsveqref}[1]{\papereqref{cpsv16}{#1}}

% Verdict marker: \gapnote{<kind>}{<status>}. Kind is the policy
% classification (clarification, local-correction, scope-restriction,
% unfaithful, false-source, open-gap); status is open, wip (elimination
% underway), resolved, or historical. Machine-read by the site index;
% prints a verdict line.
\newcommand{\gapnote}[2]{\par\noindent\textbf{Verdict:} #1 (#2).\par}


\title{Active Canonical-Form Blocks in the BNT Characterization}
\author{}
\date{2026-07-18}

\begin{document}
\maketitle

\section{The two source formulations}

Cirac--P\'erez-Garc\'ia--Schuch--Verstraete construct the canonical form by
listing the invariant blocks of the completely positive map and normalizing
each block separately \cite[lines 214--225]{Cirac2016MPDO_arXiv}. Thus the
blocks produced by the construction are precisely the blocks which occur in
the tensor. The subsequent displayed definition
\cite[lines 238--244]{Cirac2016MPDO_arXiv}, however, writes
\begin{align}
    A^i=\bigoplus_{k=1}^r \mu_k A_k^i
    \tag{\path{eq:II_CF1}}
\end{align}
without explicitly requiring every $\mu_k$ to be nonzero. Read literally, this
display also permits a listed block with zero weight. Such a block contributes
nothing to any positive-length matrix product vector.

The source definition of a basis of normal tensors
\cite[lines 271--274]{Cirac2016MPDO_arXiv} requires, as its first condition,
that every candidate representative $A_j$ is normal. Proposition~2.7 then
characterizes the basis by coverage of the canonical-form blocks and pairwise
gauge-phase distinction \cite[lines 278--280 and 1135--1146]
{Cirac2016MPDO_arXiv}. Its displayed right-hand side does not repeat candidate
normality, because that condition is already part of the term ``basis of
normal tensors.''

\section{The local correction}

To state an equivalence which is valid under the literal displayed
canonical-form hypothesis, let
\begin{align}
    K_{\mathrm{act}}=\{k:\mu_k\ne0\}.
\end{align}
Coverage is required exactly for $k\in K_{\mathrm{act}}$. Removing the other
indices leaves every positive-length matrix product vector unchanged, since
$\mu_k^N=0$ whenever $N>0$ and $\mu_k=0$. The active family has nonzero
weights and therefore satisfies the hypothesis needed by the source
phase-class argument.

Candidate normality is placed explicitly on the characterized side of the
equivalence. The resulting statement says that a candidate family is a basis
of normal tensors if and only if
\begin{enumerate}
    \item every candidate representative is normal;
    \item every active canonical-form block is gauge-phase equivalent to a
          candidate representative; and
    \item distinct candidate representatives are not gauge-phase equivalent.
\end{enumerate}
The third condition is the source notion of minimality. No converse coverage
condition is added: Definition~2.4 permits an additional normal representative
whose coefficient is identically zero, provided the full candidate family is
eventually linearly independent.

This correction is recorded in the formal characterization.\footnote{Formal
declaration:
\leanid{MPSTensor.isCPSVBasisOfNormalTensors_iff_canonicalForm_covered_and_minimal}.}
It reconciles the literal canonical-form display with the constructive
nonzero-block reading and makes the normality condition from Definition~2.4
visible in the equivalence. It does not strengthen or weaken the mathematical
characterization of the active canonical-form blocks.

\bibliographystyle{alpha}
\bibliography{references}

\end{document}
